\documentclass[12pt,a4paper]{article}

% ── Fonts and encoding (LuaLaTeX) ─────────────────────────────────────────────
\usepackage{fontspec}
\usepackage{unicode-math}
\setmainfont{Linux Libertine O}
\setmathfont{TeX Gyre Pagella Math}
\setmonofont[Scale=0.85]{Everson Mono}

% ── Layout ────────────────────────────────────────────────────────────────────
\usepackage[top=1in, bottom=1in, left=1.1in, right=1.1in, headheight=15pt]{geometry}
\usepackage{microtype}
\usepackage{parskip}
\usepackage[english]{babel}
\setlength{\emergencystretch}{3em}

% ── Headers ───────────────────────────────────────────────────────────────────
\usepackage{fancyhdr}
\pagestyle{fancy}
\fancyhf{}
\fancyhead[L]{\leftmark}
\fancyhead[R]{\thepage}
\renewcommand{\headrulewidth}{0.4pt}

% ── Section styling ───────────────────────────────────────────────────────────
\usepackage{titlesec}
\titleformat{\section}{\Large\bfseries}{\thesection}{1em}{}
\titleformat{\subsection}{\large\bfseries}{\thesubsection}{1em}{}

% ── Lists ─────────────────────────────────────────────────────────────────────
\usepackage{enumitem}
\setlist[itemize]{topsep=0.5ex, itemsep=0.25ex, parsep=0pt}
\setlist[enumerate]{topsep=0.5ex, itemsep=0.25ex, parsep=0pt}

% ── Math and tables ───────────────────────────────────────────────────────────
\usepackage{amsmath,amsthm}
\usepackage{mathtools}
\usepackage{graphicx}
\usepackage{float}
\usepackage{booktabs}
\usepackage{array}
\usepackage{tabularx}

% ── Code listings ─────────────────────────────────────────────────────────────
\usepackage{listings}
\lstset{basicstyle=\ttfamily\footnotesize, breaklines=true, columns=fullflexible}

% ── Cross-references ──────────────────────────────────────────────────────────
\usepackage{hyperref}
\usepackage{cleveref}
\hypersetup{colorlinks=true, linkcolor=blue, urlcolor=cyan, citecolor=blue,
  pdfkeywords={SIC-POVM, moduli field, conductor, ray class field, Stark units, ray class fields, Zauner conjecture, PARI/GP}}

% ── Theorem environments ──────────────────────────────────────────────────────
\newtheorem{theorem}{Theorem}[section]
\newtheorem{lemma}[theorem]{Lemma}
\newtheorem{proposition}[theorem]{Proposition}
\newtheorem{corollary}[theorem]{Corollary}
\newtheorem{definition}[theorem]{Definition}
\newtheorem{conjecture}[theorem]{Conjecture}
\newtheorem{remark}[theorem]{Remark}

% ── Custom notation ───────────────────────────────────────────────────────────
\newcommand{\Fd}{F}
\newcommand{\md}{m_d}
\newcommand{\pt}{\mathfrak{p}_2}
\newcommand{\Cl}{\mathrm{Cl}}
\newcommand{\Lean}{\textsc{Lean\,4}}
\newcommand{\Sixteen}{\textsc{sixteen}\textsubscript{3}}

\begin{document}

\title{\textbf{The Conductor of the SIC-POVM Moduli Field}}

\author{C.\ Lando\ Mills%
\thanks{The author thanks Harry T.\ Larson; see Acknowledgements and \cite{Larson1961}.}}

\date{4 August 2026}

\maketitle

\begin{abstract}
Suppose you hold the squared magnitudes of a SIC-POVM fiducial vector in
dimension~$d$. They are real, they sum to one, and their squares sum to
$2/(d+1)$. You want to know which number field they live in, exactly rather
than approximately. That field is an abelian extension of the real quadratic
field $\Fd = \mathbb{Q}(\sqrt{(d+1)(d-3)})$, and this paper determines its
conductor.

Three results carry the answer, and the third is new to this version.

At $d=12$ the moduli generate a field of degree sixteen over $\mathbb{Q}$,
signature $(8,4)$, abelian of degree eight over $\Fd$, with conductor
$\pt^3\mathfrak{p}_3$ of norm $192$ and exactly one infinite place ramified. The
field is not totally real. An earlier draft said the opposite, having been
misled by an integer relation that was fitting noise rather than recovering
structure, and the test that separates the two is whether a relation's
coefficient height holds steady as working precision rises. It is applied
throughout.

Where the class number of $\Fd$ exceeds one, the class group is not carried by
the moduli. Dimension sixteen is the smallest such dimension and it settles the
question: the $\sigma$-coinvariant count at the Appleby modulus is $16$ against
a required $d/2 = 8$, and dividing by the class group restores it. The moduli
field is the ray class field modulo the class group. At $d=2048$ this gives
degree $2^{21}$ over $\Fd$, not $2^{27}$.

Third: at $d=2048$ the rule collapses to $\pt^{12}\infty_1$, and since $2$ is
inert there, that modulus is the principal ideal generated by $4096 = 16^{3}$.
The conductor is a cube of sixteen. The paper records why that number is worth
naming rather than passing over, what supports reading it as the shape of the
sixteen-state, three-ordered trilattice this programme runs on, and what would
refute the reading.
\end{abstract}

\noindent\textbf{Keywords:} SIC-POVM \(\cdot\) moduli field \(\cdot\) conductor
\(\cdot\) ray class field \(\cdot\) Stark units \(\cdot\) archimedean ramification
\(\cdot\) Belnap trilattice \(\cdot\) Zauner conjecture \(\cdot\) PARI/GP

\newpage
\tableofcontents
\newpage

%──────────────────────────────────────────────────────────────────────────────
\section{Introduction}
\label{sec:intro}
%──────────────────────────────────────────────────────────────────────────────

Here is the problem in one sentence. A symmetric informationally complete
positive operator-valued measure in dimension~$d$ is a set of $d^2$ unit vectors
in $\mathbb{C}^d$ any two of which have the same absolute inner product,
$|\langle\psi_j,\psi_k\rangle|^2 = 1/(d+1)$ for $j \neq k$. Put less formally:
you fill complex projective space with $d^2$ points all equally far apart, and
you ask whether it can be done. Zauner conjectured in 1999 that it can, in every
dimension, and that the vectors can be taken as the orbit of a single fiducial
$\psi$ under the $d^2$ Weyl--Heisenberg displacement operators~\cite{Zauner1999}.
The conjecture is open.

Here is why it is a number theory problem. Appleby, Bengtsson, Flammia,
McConnell, and Yard~\cite{Appleby2013,Appleby2017,AFMY2020} established that the
fiducial coordinates are algebraic numbers living in ray class fields of the
real quadratic field $\Fd = \mathbb{Q}(\sqrt{\md})$, with $\md$ the squarefree
part of $(d+1)(d-3)$. That ties SIC existence to the real-quadratic case of
Hilbert's twelfth problem. Find the fiducial and you have
a generator; find the conductor and you know which field to look in.

That is the gap this paper fills. Knowing a fiducial lives in \emph{some} ray
class field is not knowing \emph{which}. For any explicit construction the
conductor is the thing you need, because it determines the ray class group, the
number of characters, and whether the Stark machinery applies at all. This paper
determines the conductor for the moduli, the real quantities $N_k = |\psi_k|^2$.
The determination is empirical in the honest sense: calibrated against
dimensions where the fiducial is known exactly, then applied to one where it is
not.

The rule is stated in full in \Cref{sec:rule}. Writing $\pt$ for the prime of
$\Fd$ above $2$, the moduli generate the ray class field of $\Fd$ modulo the
class group, at the modulus
\[
\mathfrak{f}_d \;=\; \pt^{\,v_2(d)+1} \cdot \prod_{p \mid d,\; p \text{ odd}}
\ \prod_{\mathfrak{p} \mid p} \mathfrak{p} \cdot \infty_1,
\]
with exactly one infinite place ramified.

Three features of this answer are more important than the formula.

The first is archimedean. Exactly one infinite place ramifies, so the moduli
field is real at the fiducial embedding but not totally real. Stark's machinery
in its standard form wants a totally real class field and refuses anything else;
PARI's \texttt{bnrstark} reports nonzero $r_2$ and exits. The construction needs
the mixed-signature form. Concretely, a Stark unit here is recovered as a monomial
in $S$-units of $\Fd$ rather than from a totally real Stark--Brumer construction,
which is what \Cref{sec:unit} carries out at the predicted modulus.

The second is the class group, invisible in every dimension where the fiducial
is known exactly because all of those have class number one, and settled at
$d=16$ in \Cref{sec:classgroup}. It costs six powers of two at $d=2048$.

The third is what the rule produces at $d=2048$, and it is the reason this
version exists. There $d = 2^{11}$ has no odd prime divisor, so the formula
collapses to $\pt^{12}\infty_1$; and $2$ is inert in $\Fd$, so $\pt = (2)$ and
the modulus is the principal ideal
\[
  \mathfrak{f}_{2048} \;=\; (2^{12})\,\infty_1 \;=\; (4096)\,\infty_1
  \;=\; (16^{3})\,\infty_1 .
\]
The conductor is generated by sixteen cubed. \Cref{sec:sixteen} treats that as a
fact worth naming rather than an arithmetic accident, and is explicit about
which part of the treatment is arithmetic and which is identification.

%──────────────────────────────────────────────────────────────────────────────
\section{The moduli field is not the coordinate field}
\label{sec:notcoord}
%──────────────────────────────────────────────────────────────────────────────

It is tempting to treat the moduli field and the coordinate field as
interchangeable. They are not. For a fiducial $\psi$ the coordinates $\psi_k$
are complex; the moduli $N_k = |\psi_k|^2$ are real. But reality at the fiducial
embedding does not imply total reality: the field can be, and is, complex at the
conjugate place.

The practical teeth are that PARI's \texttt{bnrstark} computes Stark units for
totally real class fields and refuses anything else. It refuses the moduli
field. An earlier version of this paper claimed the opposite on the strength of
a computation that did not survive scrutiny. Here is what survived.

\begin{proposition}[$d=12$ moduli field]
\label{prop:d12}
The twelve moduli of the $d=12$ fiducial generate a field of degree sixteen over
$\mathbb{Q}$ of signature $(8,4)$. Individually they have degree $4$ or $8$; the
minimal polynomial of $N_0$ is
$97344\,x^4 - 32448\,x^3 + 3224\,x^2 - 104\,x + 1$. The field is abelian of
degree eight over $\Fd = \mathbb{Q}(\sqrt{13})$, with conductor
$\pt^3\mathfrak{p}_3$ of norm $192$, \emph{one} infinite place ramified, ray
class group $(\mathbb{Z}/2)^3$ of order eight, and trivial subgroup.
\end{proposition}

The field is therefore not totally real, and the moduli field is the full ray
class field at a modulus whose archimedean part is one place: not none, and not
both.

Establishing this takes more care than it appears, and the care is the real
content. An integer relation computed against a power basis will happily report
membership in a field that does not contain the element, if the working
precision gives the lattice reduction enough slack to fit noise. The
discriminating test is not the residual, since small residuals are easy. The
test is the height. A genuine algebraic relation keeps its coefficient height
fixed as working precision rises; a fitted relation sees its height grow with
the precision, because each new digit demands a new coefficient to match it.

At $d=12$, $N_1$ has a degree-eight minimal polynomial of height $33$ bits,
unchanged at $600$, $900$, and $1300$ digits. The primitive element
$\sum_k k N_k$ has a degree-sixteen relation of height $80$ bits, likewise
unchanged. But the claim that $N_1$ lies in a totally real octic, which is the
claim the earlier draft made, is returned at every one of the eight embeddings
with heights of $177$, $287$, and $435$ bits at those three precisions. That is
the signature of a fit, not a fact. The moduli have degree at most eight each,
they do not lie in a common octic, and together they generate the degree-sixteen
field above.

A related trap governs the input. The exact coordinates are known to roughly
sixteen hundred digits, so seeking a relation at twenty-five hundred is fitting
digits that are not there. Every claim in this paper is stated at a precision
the input supports.

%──────────────────────────────────────────────────────────────────────────────
\section{How each field was identified}
\label{sec:method}
%──────────────────────────────────────────────────────────────────────────────

The method is the same at each calibration dimension and is stated once.

Take the fiducial to the precision the published data supports. Form the moduli
$N_k = |\psi_k|^2$. Seek a minimal polynomial for each by integer relation
against a power basis, and seek one for a primitive element, taken here as
$\sum_k k N_k$. Accept a relation only if its coefficient height is stable
across at least three working precisions spanning several hundred digits.
Identify the field from the accepted polynomial, compute its signature, and
verify abelian-ness over $\Fd$ by checking that the degree divides the ray class
group order at a candidate modulus and that the Galois group is what the ray
class group says it should be.

The conductor is then read off rather than guessed: enumerate moduli of $\Fd$ of
small norm, compute the ray class group at each with the infinite places treated
both ways, and select the one whose group has the right order and type and whose
subgroup is trivial. Where several moduli agree the smallest is taken, which is
what conductor means.

Two warnings, both learned rather than anticipated. The height test above is
necessary and not decorative. And the archimedean part must be searched, not
assumed: taking both infinite places unramified because the moduli are real is
exactly the error that produced the retracted claim.

%──────────────────────────────────────────────────────────────────────────────
\section{The three calibration points}
\label{sec:calibration}
%──────────────────────────────────────────────────────────────────────────────

Three dimensions have fiducials known exactly and class number one: $d=4$,
$d=8$, $d=12$. All three were computed by the method of \Cref{sec:method}.

Dimension four has $v_2(4) = 2$ and receives conductor exponent three.
Dimension eight has $v_2(8) = 3$ and receives exponent four. Dimension twelve
has $v_2(12) = 2$ and receives exponent three, together with the prime above
three, since three divides twelve.

The reason dimension eight was worth the effort is that it refutes the naive
guess. A constant exponent of three fits dimensions four and twelve perfectly
well. Only $d=8$ separates the constant from $v_2(d)+1$, and it does so
decisively.

The odd part carries a subtlety worth appreciating. Three splits in
$\mathbb{Q}(\sqrt{13})$, so the factor appearing in the $d=12$ conductor has
norm three, while three is inert in $\mathbb{Q}(\sqrt5)$ with norm nine. The odd
contribution is a statement about primes of $\Fd$ lying above primes dividing
$d$, not about rational integers. The two cases would be indistinguishable if
only one had ever been computed, which is another way of saying that three
calibration points is the minimum rather than a luxury.

%──────────────────────────────────────────────────────────────────────────────
\section{The rule}
\label{sec:rule}
%──────────────────────────────────────────────────────────────────────────────

\begin{conjecture}[Conductor of the moduli field]
\label{conj:rule}
Let $d$ be even, $\Fd = \mathbb{Q}(\sqrt{\md})$ with $\md$ the squarefree part
of $(d+1)(d-3)$, and $\pt$ the prime of $\Fd$ above $2$. The moduli
$N_k = |\psi_k|^2$ of a Zauner-symmetric Weyl--Heisenberg fiducial in dimension
$d$ generate the ray class field of $\Fd$ at the modulus
\[
\mathfrak{f}_d \;=\; \pt^{\,v_2(d)+1}\cdot \prod_{p\mid d,\ p\ \mathrm{odd}}\
\prod_{\mathfrak{p}\mid p}\mathfrak{p}\cdot\infty_1,
\]
with exactly one infinite place ramified, taken modulo the class group of $\Fd$.
Its degree over $\Fd$ is $|\Cl_{\mathfrak{f}_d}|/h(\Fd)$. The field is real at
the fiducial embedding and not totally real.
\end{conjecture}

The formula is compact and every piece earns its place. The $2$-adic exponent is
forced by $d=8$ against $d=4$ and $d=12$. The odd part is forced by the splitting
behaviour just described. The archimedean part is forced by
\Cref{prop:d12}. The quotient by the class group is forced by $d=16$, which is
the subject of \Cref{sec:classgroup}.

%──────────────────────────────────────────────────────────────────────────────
\section{Dimension 2048}
\label{sec:d2048}
%──────────────────────────────────────────────────────────────────────────────

Two thousand and forty-eight is the eleventh power of two. It has no odd prime
divisors, so \Cref{conj:rule} collapses to $\mathfrak{f}_{2048} = \pt^{12}\infty_1$.
Here $\md = 4190205 = 3\cdot 5\cdot 409\cdot 683$, the class number is $64$ with
class group $[32,2]$, and two is inert in $\Fd$ since $\md \equiv 5 \pmod 8$.

Because two is inert, $\pt = (2)$, and the modulus is principal:
\begin{equation}\label{eq:modulus}
  \mathfrak{f}_{2048} \;=\; (2^{12})\,\infty_1 \;=\; (4096)\,\infty_1 .
\end{equation}

The ray class group at $\pt^{12}\infty_1$ has order $2^{27}$, with cyclic type
$[4096,1024,8,4]$. The class number is $64$, so by \Cref{sec:classgroup} the
moduli field has degree $2^{21}$ over $\Fd$ and $2^{22}$ over $\mathbb{Q}$.

To see that this prediction is not sitting on a plateau where an error of one
would be invisible, the neighbouring exponents were computed with the infinite
places unramified, so that growth in the finite part stands on its own.

\begin{center}
\begin{tabular}{@{}clll@{}}
\toprule
exponent $k$ & $|\Cl_{\pt^k}|$ & $\log_2$ & cyclic type \\
\midrule
$9$  & $2097152$   & $21$ & $[1024, 128, 8, 2]$ \\
$10$ & $8388608$   & $23$ & $[2048, 256, 8, 2]$ \\
$11$ & $33554432$  & $25$ & $[4096, 512, 8, 2]$ \\
$12$ & $67108864$  & $26$ & $[4096, 1024, 8, 2]$ \\
$13$ & $134217728$ & $27$ & $[4096, 2048, 8, 2]$ \\
\bottomrule
\end{tabular}
\end{center}

The growth is a factor of four per step up to the eleventh power and a factor of
two from there. The predicted exponent falls exactly at the change of régime. An
error of one either way is an error of a factor of two or four in the degree,
which is not subtle.

Read in valuations the filtration has a single generating shape, and
\Cref{fig:filtration} shows it. Writing each cyclic invariant as its $2$-adic
valuation, every level from $k=5$ upward is $[a,\,b,\,3,\,1]$: the tail is frozen
and never moves again. The gap between the two leading invariants is $a-b=3$,
equal to the valuation of the third invariant, and it holds for the whole
fourfold r\'egime. Growth of four is therefore not two independent doublings but
one invariant advancing while its distance to the next is pinned at the tail's
valuation.

The geometry of that statement is a staircase with a ceiling. The leading
invariant climbs one valuation per level, which is the ray class group acquiring
one factor of two per step in the finite part of the modulus; the second
invariant climbs in parallel three valuations behind, so the pair advances as a
rigid body and the group grows by four. Nothing in the tail participates: from
$k=5$ the third and fourth invariants sit at $3$ and $1$ and never move, so the
whole of the growth lives in a two-dimensional slice of a four-dimensional
filtration. The r\'egime changes exactly when the leader saturates at
$12 = v_2(2d)$, which is a ceiling and not a coincidence: the leading invariant
cannot exceed the $2$-adic content of $2d$, so once it has absorbed all of it,
only the second can still move. The rigid body stops being rigid. The gap
collapses through $2,1,0$ at $k=12,13,14$, the second invariant closes the
distance and catches the first at $k=14$, and at $k=15$ the leader resumes.
Growth is two throughout, because exactly one invariant is moving at a time.

The predicted exponent is the last level at which the gap is still pinned, the
last moment before the structure changes character. It is therefore not a point
chosen inside a smooth r\'egime, where an error of one would be invisible; it is
the corner itself.

\begin{figure}[H]
\centering
\includegraphics[width=\linewidth]{figs/filtration_2048.pdf}
\caption{The $2$-adic filtration of the moduli field at $d=2048$, computed with
PARI/GP and closed at $k=15$. Above: the valuation of each cyclic invariant of
$\Cl_{\pt^{k}}$. The tail is frozen from $k=5$; the leading pair advances as a
rigid body until the leader meets the ceiling $v_2(2d)=12$. Below: the gap
$a-b$ between the two leading valuations, pinned at the tail's valuation
throughout the fourfold r\'egime and collapsing through $2,1,0$ once the ceiling
is reached. The predicted exponent $k=12$ is the corner.}
\label{fig:filtration}
\end{figure}

The calibration and the prediction are recorded in \Lean{} alongside the tower
data, with the exponent rule and the fact that dimensions four and eight share a
base field but differ in exponent stated as decidable propositions and
discharged by evaluation. The tower has since been carried
to $k=15$, which closes it: every ray class degree in the \Lean{} development is
computed rather than extrapolated, and the halving is sustained, holding at
exactly two for each of $k=12,13,14,15$.

%──────────────────────────────────────────────────────────────────────────────
\section{The unit at the conductor}
\label{sec:unit}
%──────────────────────────────────────────────────────────────────────────────

The conductor names a field. Before any fiducial is recovered there are
elements of $\Fd$ that the conductor makes available, and they are available in
closed form at every even $d$ rather than only at $d=2048$. This section records
them, because the arithmetic is uniform and the uniformity is what makes it
usable.

The starting point is an identity about $\md$ itself. Since
$\md = (d+1)(d-3)$,
\[
  \md \;=\; d^2 - 2d - 3 \;=\; (d-1)^2 - 4 .
\]
Two consequences follow immediately.

\begin{proposition}[The unit and the two ramified generators]
\label{prop:unit}
Let $d$ be even and $\md = (d+1)(d-3)$ be the squarefree part as above. Put
\[
  \varepsilon_d = \frac{(d-1) + \sqrt{\md}}{2}, \qquad
  g_- = \frac{\sqrt{\md} - (d-3)}{2}, \qquad
  g_+ = \frac{(d+1) - \sqrt{\md}}{2}.
\]
Then $\varepsilon_d$ satisfies $x^2 - (d-1)x + 1 = 0$ and has norm $1$, so it is
a unit of $\Fd$ whose Galois conjugate is its inverse; and
\[
  N(g_-) = -(d-3), \qquad N(g_+) = d+1 .
\]
Writing $\sigma(\varepsilon_d) = \varepsilon_d^{-1}$ for the conjugate
embedding, the three are not independent:
\[
  g_- = 1 - \sigma(\varepsilon_d), \qquad g_+ = 1 + \sigma(\varepsilon_d).
\]
\end{proposition}

\begin{proof}
$N(\varepsilon_d) = \bigl((d-1)^2 - \md\bigr)/4 = 4/4 = 1$ by the identity
above. For $g_-$, $\bigl((d-3)^2 - \md\bigr)/4 = (d-3)\bigl((d-3)-(d+1)\bigr)/4
= -(d-3)$; for $g_+$, $\bigl((d+1)^2 - \md\bigr)/4 = (d+1)\bigl((d+1)-(d-3)\bigr)/4
= d+1$. The last two identities are direct substitution of
$\sigma(\varepsilon_d) = \bigl((d-1) - \sqrt{\md}\bigr)/2$.
\end{proof}

So the two elements whose norms are the two rational factors of $\md$ are the
unit's conjugate displaced by one in either direction. That is a fact about the
shape of $\md$, and it holds wherever the conductor rule of \Cref{conj:rule}
applies.

\subsection{The $S$-unit monomial at exponents $[-1,3,2]$}

The monomial that the conductor-$16$ tower selects is
$M_d = \varepsilon_d^{-1} g_-^{\,3} g_+^{\,2}$. By \Cref{prop:unit} it is a
polynomial in the single quantity $\sigma = \sigma(\varepsilon_d)$:
\begin{equation}\label{eq:monomial-sigma}
  M_d \;=\; \sigma\,(1-\sigma)^3(1+\sigma)^2
        \;=\; \sigma - \sigma^2 - 2\sigma^3 + 2\sigma^4 + \sigma^5 - \sigma^6 .
\end{equation}
No division occurs anywhere: $\varepsilon_d$ has norm one, so its inverse is its
conjugate, and the monomial is a product.

\begin{proposition}[What the monomial approximates]
\label{prop:monomial}
$M_d = \dfrac{1}{d} - \dfrac{2}{d^{3}} + O\!\left(d^{-4}\right)$.
\end{proposition}

\begin{proof}
Write $N = d-1$. From $\sigma = \bigl(N - \sqrt{N^2-4}\bigr)/2$ one has
$\sigma = N^{-1} + N^{-3} + O(N^{-5})$. Substituting into
\Cref{eq:monomial-sigma} gives $M_d = N^{-1} - N^{-2} - N^{-3} + O(N^{-4})$,
while $1/d = 1/(N+1) = N^{-1} - N^{-2} + N^{-3} - O(N^{-4})$. Subtracting,
$M_d - 1/d = -2N^{-3} + O(N^{-4})$.
\end{proof}

The exponent vector $[-1,3,2]$ therefore names $1/d$ to third order in every
dimension, not only at $d=2048$. Computation confirms the constant and shows the
approach:

\begin{center}
\begin{tabular}{@{}rll@{}}
\toprule
$d$ & $M_d - 1/d$ & $\bigl(M_d - 1/d\bigr)\,d^{3}$ \\
\midrule
$16$      & $-5.78\times10^{-4}$  & $-2.367$ \\
$256$     & $-1.20\times10^{-7}$  & $-2.020$ \\
$2048$    & $-2.33\times10^{-10}$ & $-2.002$ \\
$4096$    & $-2.91\times10^{-11}$ & $-2.001$ \\
$10^{7}$  & $-2.00\times10^{-21}$ & $-2.000001$ \\
\bottomrule
\end{tabular}
\end{center}

That the law is dimension-independent is the useful half of the statement. A
monomial that named $1/d$ only at $d=2048$ would be a coincidence of size; one
that names it everywhere is a property of the tower, and it is what allows the
representation below to be carried between dimensions.

At $d=2048$ the monomial is exactly
\begin{equation}\label{eq:monomial-2048}
  M_{2048} \;=\;
  \frac{-73534985899138338825 \;+\; 35923312651511805\,\sqrt{4190205}}{2},
\end{equation}
with norm $N(M_{2048}) = 1\cdot(-2045)^3\cdot 2049^2 = -35905737691441125$,
obtained multiplicatively from \Cref{prop:unit} rather than from
$(a^2 - b^2\md)/c^2$, which exceeds any fixed machine width well before it is
needed.

\subsection{Exactness is not evaluability}

\Cref{eq:monomial-2048} is exact and is also the one form of that number which
double precision cannot read. Its two terms are each near $7.35\times10^{19}$
and agree through twenty significant digits; the value that survives the
subtraction is near $4.9\times10^{-4}$. The unit in the last place at that
magnitude is larger than the answer by twenty-three orders, so evaluating the
difference returns rounding and nothing else.

The cancellation is in the reading, not in the number, and rationalizing removes
it. For $M = (a + b\sqrt{\md})/c$,
\begin{equation}\label{eq:rationalize}
  M \;=\; \frac{c\,N(M)}{\,a - b\sqrt{\md}\,},
\end{equation}
where the denominator now \emph{adds} two terms of like sign, so no significance
is lost, and the numerator is an exact rational integer available
multiplicatively. Evaluated this way \Cref{eq:monomial-2048} gives
$0.000488281016884806$, correct in every digit shown.

This is the numerical-analytic counterpart of the height test of
\Cref{sec:notcoord}, and it deserves the same standing. There the caution is
that a relation can be reported at a precision the input does not support; here
it is that an exact representation can be unreadable in the arithmetic used to
read it. Neither failure announces itself: the fitted relation has a small
residual, and the cancelling expansion returns a number of ordinary size.

\subsection{The radicals from their cyclotomic levels}

$\sqrt{\md}$ at $d=2048$ is available from the levels its factorization names.
With $g(m) = \sum_{k=0}^{m-1} e^{2\pi i k^2/m}$ the quadratic Gauss sum, the
radical is read from a component of the sum and not from its modulus: at a prime
level $p \equiv 1 \pmod 4$ the sum is real and equals $\sqrt{p}$, while at a
level divisible by four it is $2\sqrt{n}\,(1+i)$, whose modulus is $\sqrt{8n}$
and whose real part halved is $\sqrt{n}$.

\begin{center}
\begin{tabular}{@{}rlll@{}}
\toprule
level $m$ & $g(m)$ & radical & extracted from \\
\midrule
$5$    & $2.236067977$                & $\sqrt{5} = 2.236067977$    & $\operatorname{Re}$ \\
$12$   & $3.464101615(1+i)$           & $\sqrt{3} = 1.732050808$    & $\operatorname{Re}/2$ \\
$409$  & $20.223748416$               & $\sqrt{409} = 20.223748416$ & $\operatorname{Re}$ \\
$2732$ & $52.268537381(1+i)$          & $\sqrt{683} = 26.134268691$ & $\operatorname{Re}/2$ \\
\bottomrule
\end{tabular}
\end{center}

Each radical lands within $1.5\times10^{-13}$ of its target, their product
recovers $\sqrt{\md} = 2046.999022960$ to $1.3\times10^{-11}$, and
$\varepsilon_{2048} = \bigl(2047 + \prod\bigr)/2$ follows from the sums rather
than from a seeded constant. Comparing $|g(12)|$ with $\sqrt3$ instead compares
$\sqrt{24}$ with $\sqrt3$; the modulus is simply the wrong component, not a
failed identity.

\subsection{What the representation costs}

The practical claim attached to an exact representation is that it is small.
Counting the actual widths of \Cref{eq:monomial-2048} across dimensions makes
the claim quantitative. Writing $M_d = (a+b\sqrt{\md})/c$, the bit-lengths are
$\lceil\log_2 a\rceil = 6\log_2 d$, $\lceil\log_2 b\rceil = 5\log_2 d$, and
$\lceil\log_2 \md\rceil = 2\log_2 d$, so the whole datum costs about
$13\log_2 d$ bits:

\begin{center}
\begin{tabular}{@{}rrrr@{}}
\toprule
$d$ & algebraic (bits) & $\div \log_2 d$ & $d^2$ amplitudes (bits) \\
\midrule
$16$      & $54$  & $13.50$ & $32768$ \\
$256$     & $106$ & $13.25$ & $8388608$ \\
$2048$    & $145$ & $13.18$ & $536870912$ \\
$65536$   & $210$ & $13.12$ & $549755813888$ \\
$2^{20}$  & $262$ & $13.10$ & $1.41\times10^{14}$ \\
\bottomrule
\end{tabular}
\end{center}

The ratio is still falling toward its asymptote at $2^{20}$, which is the shape
a logarithmic law has. What is being compared is a representation against a
representation: this is the cost of naming the algebraic datum, not the cost of
producing a fiducial, and \Cref{sec:follows} is explicit about which of those
$d=2048$ still lacks.

%──────────────────────────────────────────────────────────────────────────────
\section{The modulus is sixteen cubed}
\label{sec:sixteen}
%──────────────────────────────────────────────────────────────────────────────

\Cref{eq:modulus} says the conductor at $d=2048$ is generated by $4096$. That
number is $2^{12}$, and it is also $16^{3}$. This section says why the second
reading is worth recording.

The programme this paper belongs to carries a logical substrate that is neither
Boolean nor merely four-valued. Its base is the four-element set
\[
  I = \{\,T,\ F,\ t,\ f\,\}
\]
of constructively proven, constructively refuted, non-constructively acceptable,
and non-constructively rejectable. The trilattice is the full powerset
$\mathcal{P}(I)$, so it has $2^{4} = 16$ states, and it carries three genuinely
distinct orderings:
\begin{align}
  x \le_{i} y &\iff x \subseteq y, \label{eq:order-i}\\
  x \le_{t} y &\iff x \cap \{T,t\} \subseteq y \cap \{T,t\}
                \ \text{ and }\ y \cap \{F,f\} \subseteq x \cap \{F,f\},
                \label{eq:order-t}\\
  x \le_{c} y &\iff x \cap \{T,F\} \subseteq y \cap \{T,F\}
                \ \text{ and }\ y \cap \{t,f\} \subseteq x \cap \{t,f\}.
                \label{eq:order-c}
\end{align}
Information, truth, constructivity: sixteen states, three orders
\cite{Belnap1977,Shramko2005}. The pair $(16,3)$ travels with the structure.

The passage from the four to the sixteen deserves a statement, because it
happens twice over and for a reason that does not generalise.

\begin{proposition}\label{prop:coincidence}
$|\mathcal{P}(I)| = 2^{4} = 16$ and $|I|^{2} = 4^{2} = 16$. The equation
$2^{n} = n^{2}$ has exactly two solutions in the positive integers, $n = 2$ and
$n = 4$.
\end{proposition}

\begin{proof}
Check $n = 1,2,3,4$ directly: $2 \neq 1$, $4 = 4$, $8 \neq 9$, $16 = 16$. For
$n \ge 5$ argue by induction that $2^{n} > n^{2}$. The base case is
$2^{5} = 32 > 25$. If $2^{n} > n^{2}$ with $n \ge 5$ then
$2^{n+1} = 2\cdot 2^{n} > 2n^{2} \ge (n+1)^{2}$, the last step being equivalent
to $n^{2} - 2n - 1 \ge 0$, which holds for $n \ge 3$.
\end{proof}

So four is the only place above the trivial case where enumerating all subsets
and letting the base refer to itself give the same count. That is what licenses
reading sixteen as \emph{four self-applied} rather than merely \emph{four
enumerated}, and nothing more general licenses it.

\begin{conjecture}\label{conj:identification}
The $16$ and the $3$ in $\mathfrak{f}_{2048} = (16^{3})\infty_1$ are the
trilattice's sixteen states and three orderings: the conductor of the $d=2048$
moduli field is the trilattice raised to its own number of orders.
\end{conjecture}

Three independent things support this, and none of them is inflated by saying so.

\begin{figure}[H]
\centering
\includegraphics[width=\linewidth]{figs/horn_side.pdf}
\caption{The side-on section, in the plane of the syzygy axis, drawn by the
same generator that produces \Cref{fig:horn}. The $\in$ shell is centred on
$\odot$; the evaluator sphere carries the same radius and is centred one radius
along the axis, so $\odot$ lies exactly on it and the two are tangent. The
segment between the centres is the coupler. The evaluator trine is seen edge-on
in this plane and stands at $120^{\circ}$ only in the perpendicular one, which
is why the two sections are drawn separately rather than combined.}
\label{fig:cfg}
\end{figure}

\textbf{The arithmetic is exact.} \Cref{eq:modulus} is not an approximation and
contains no tuned parameter. Either $\pt^{12}$ is principal and generated by
$4096$ or it is not, and inertness of $2$ settles that it is.

\textbf{The same pair is already present geometrically.} Independently of any
ray class computation, the evaluator arms of this programme trace a $(16,3)$
torus knot on the horn torus, sixteen toroidal windings against three poloidal
ones, with bridge number three matching the three arms that carry the
four-valued kernel. That pair was in the programme before this conductor was
computed, and \Cref{fig:cfg} shows that pair of spheres in section.

\Cref{fig:horn} shows the equator that claim is made on. The horn torus is the
degenerate case $R=r$, where the tube radius equals the centre radius and the
hole closes to a single point. That point is the pinch, and it is not a defect
of the drawing: it is where the boundary meets itself, which is why the surface
can carry a self-referential reading at all. The three evaluator arms sit at
$120^{\circ}$ on the equator, the $A_2$ geometry, and the three poloidal
windings of the $(16,3)$ knot are those same three arms seen as a circuit rather
than as a triple. The sixteen toroidal windings are the trilattice's sixteen
states carried once around. Nothing here is scaled for legibility: $R=r$ exactly,
so the two lobes of the side-on view touch rather than overlap, which is the
whole content of the pinch.

\begin{figure}[H]
\centering
\includegraphics[width=\linewidth]{figs/horn_equator.pdf}
\caption{The equatorial section of the horn torus, drawn at $R = r$. The pinch
$\odot$ sits at the origin, where the hole has closed to a point and the
boundary meets itself. The three evaluator arms stand at $120^{\circ}$ and span
the $A_2$ triangle; they are the three poloidal windings of the $(16,3)$ knot,
seen as a circuit rather than as a triple. Rendered from the generator that
draws the kernel geometry, in its print theme.}
\label{fig:horn}
\end{figure}

\textbf{The exponent and the calibrating dimension coincide.} The leading cyclic
invariant saturates at $12 = v_2(2d)$, which is also the conductor exponent and
also the dimension at which the four-valued base is carried on twelve axes. We
note the coincidence and deliberately do not build on it; three occurrences of
the integer twelve in one programme is the kind of thing that wants a reason
before it wants a use.

\begin{remark}[What would refute it]
The reading fails if the conductor at other $d = 2^{n}$ is not a power of
sixteen. The rule gives $\mathfrak{f}_{2^{n}} = \pt^{n+1}\infty_1$, principal
whenever $2$ is inert, so the generator is $2^{n+1}$ and is a power of sixteen
exactly when $4 \mid n+1$. That happens at $n = 3, 7, 11, 15$, giving
$d = 8, 128, 2048, 32768$ with generators $16, 16^{2}, 16^{3}, 16^{4}$. So the
claim is not that every dimension has a sixteen-power conductor; it is that the
sixteen-power dimensions form the arithmetic progression $n \equiv 3 \pmod 4$,
and that $d = 2048$ is the third of them. If the $d=8$ conductor were not
generated by $16$, or if the $(16,3)$ knot invariance failed somewhere in that
progression, or if the three orderings of
\Cref{eq:order-i,eq:order-t,eq:order-c} proved dependent so that the structure
is a bilattice, the reading should be dropped.
\end{remark}

\begin{remark}
Dimension eight is the first member of that progression, and it is one of the
three calibration points where the fiducial is known exactly. Its conductor is
$\pt^{4}$, and $\md = 5$ there with $2$ inert, so the generator is $16$ itself.
The progression therefore begins at a dimension where the answer is not
predicted but computed, which is the strongest position the claim can be in
without the family-wide computation.
\end{remark}

%──────────────────────────────────────────────────────────────────────────────
\section{The class group, settled at dimension sixteen}
\label{sec:classgroup}
%──────────────────────────────────────────────────────────────────────────────

Every ray class field contains the Hilbert class field. So a rule naming the
full ray class field and a rule naming its quotient by the class group differ by
a factor of $h(\Fd)$. All three calibration dimensions have $h=1$, where the two
agree, so nothing in \Cref{sec:calibration} tells them apart. To see the class
group you need a dimension where it is nontrivial, and dimension sixteen is the
smallest.

At $d=16$ we have $\md = 221 = 13\cdot 17$, class number two, two inert in
$\Fd = \mathbb{Q}(\sqrt{221})$, and $v_2(16)+1 = 5$. The modulus is
$\pt^5\infty_1$, the ray class group has order $256$ with cyclic type $[16,8,2]$,
giving $128$ over $\Fd$ after the class group is divided out. The $2$-power
tower below is listed with the infinite places unramified, so that growth in the
finite part is visible on its own.

\begin{center}
\begin{tabular}{@{}clll@{}}
\toprule
$k$ & $|\Cl_{\pt^k}|$ & cyclic type & narrow order \\
\midrule
$3$ & $16$  & $[8,2]$   & $64$ \\
$4$ & $64$  & $[16,4]$  & $256$ \\
$5$ & $128$ & $[16,8]$  & $512$ \\
$6$ & $256$ & $[16,16]$ & $1024$ \\
\bottomrule
\end{tabular}
\end{center}

But the deciding quantity is not the order of this group. It is the count of
independent moduli the group has to carry. Galois conjugation by the nontrivial
automorphism $\sigma$ of $\Fd$ pairs the moduli, $N_{k+d/2} = \sigma(N_k)$, so a
moduli field must support exactly $d/2$ of them, and that count is read off the
$\sigma$-coinvariant quotient of the ray class group at the Appleby modulus
$(3d)$.

\begin{proposition}[The coinvariant count at $d=16$]
\label{prop:coinv}
Write $G_d$ for the ray class group of $\Fd$ at $(3d)$ with both infinite places
unramified, and $G_d^\sigma$ for its $\sigma$-coinvariants. Then
$|G_{16}^\sigma| = 16$, which is not $d/2 = 8$, while
\[
\frac{|G_{16}^\sigma|}{|\Cl(\Fd)^\sigma|} \;=\; \frac{16}{2} \;=\; 8 \;=\; \frac{d}{2}.
\]
The same identity holds at $d = 4,\,8,\,12$, where the class number is one and
the denominator is trivial, so those dimensions say nothing either way. It is
$d=16$ that carries the content.
\end{proposition}

This is not a law for every dimension, and the reason is arithmetic rather than
accidental, so it is worth setting down. It holds at $d = 4, 8, 12, 16, 24, 32,
36$ and fails at $d = 20, 28, 40$. Quoting the corrected count throughout, the
coinvariant order divided by $|\Cl(\Fd)^\sigma|$, the three failures give $8$,
$12$, and $16$ against $d/2$ of $10$, $14$, and $20$.

The dividing line is the odd part of $d/2$. For an odd prime $q$, a ray class
group at a modulus divisible by $\mathfrak{q}$ but not $\mathfrak{q}^2$ has no
$q$-torsion: the local unit group has order $N\mathfrak{q}-1$, equal to $q-1$
when $q$ splits and $q^2-1$ when $q$ is inert, and $q$ divides neither. The
$q$-torsion first appears at $\mathfrak{q}^2$. The modulus $(3d)$ supplies $3^2$
exactly when $3 \mid d$, and supplies no other odd prime squared unless that
square already divides $d$. So the count reaches $d/2$ when the odd part of
$d/2$ is a power of three, and falls short otherwise. Supplying the square by
hand does not repair it: at $d=20$ the moduli $25$ and $100$ give counts $20$
and $40$ where $10$ is wanted; at $d=28$ the modulus $49$ gives $21$ where $14$
is wanted.

So \Cref{prop:coinv} is a statement about dimension sixteen and not an instance
of a general law, and it should be weighed as one dimension. What it says there
is nonetheless exact: the raw count overshoots by precisely the class number,
and dividing by the class group restores the number of independent moduli that
Galois pairing requires.

Together with the typing of the two readings, where the reading that keeps the
class group grounds itself on a non-abelian Galois group that the ray class
group $\mathbb{Z}/16\times\mathbb{Z}/8$ does not have, the conclusion is clean.
The class group is not carried by the moduli. The moduli field is the ray class
field at $\mathfrak{f}_d$ modulo the class group, of degree
$|\Cl_{\mathfrak{f}_d}|/h(\Fd)$ over $\Fd$: degree $128$ at $d=16$, and
$2^{27}/64 = 2^{21}$ at $d=2048$. The calibration of \Cref{sec:calibration} is
untouched, and what changes is only the reading of the word \emph{full}. It
changes the $d=2048$ degree by six powers of two.

One further fact about $d=16$ deserves a place in the record. The $d=16$ SIC is
in hand numerically: the Zauner eigenspaces have dimensions five, six, and five;
only the six-dimensional sector carries fiducials; and eight of them, separated
by their moduli multisets, refine against the full overlap system past two
thousand digits. The sum of moduli and the sum of their squares are exact to
precision.

That computation also fixes the limit of recognition from below. With the moduli
field of degree $256$ over $\mathbb{Q}$ at $d=16$, an individual modulus has no
minimal polynomial within reach of integer relation. Indeed $\sqrt{221}$, which
certainly lies in the field, is not recovered from the power basis of a modulus
to degree twenty at twenty-five hundred digits. The height-collapse signature
that identifies the canonical fiducial at $d=8$, where the field has degree
eight, does not survive the growth of the field. The route at $d=16$ and beyond
is from the field down, not from the numbers up.

%──────────────────────────────────────────────────────────────────────────────
\section{The tower data, and what asserting it concealed}
\label{sec:discharge}
%──────────────────────────────────────────────────────────────────────────────

The tower entered \Lean{} as thirty axioms: the degree at each of sixteen
levels, the $2$-adic valuations, the leading and trailing cyclic invariants, the
class number, the narrow degree, and three field types. That is an honest way to
record computed field data, and it is what the artifact statement means by the
boundary between what is proved and what is assumed. It is also more than the
data requires.

The sixteen degrees are not independent. The same module states the two laws
they obey, a ratio of four through the maximal-growth interval $3\le k\le11$ and
of two once the leading invariant saturates, so only the three seed values $64$,
$64$, $128$ are data. Written as a recursion on those, every level is decidable
and the growth law becomes a consequence rather than a companion assumption. The
valuations follow, though not by decision: $\nu_2$ is \texttt{padicValNat}, which
is noncomputable, and each level is instead an instance of $\nu_p(p^n)=n$,
available because every degree in the tower is a power of two. The narrow degree
is four times the wide one, the factor this manuscript already attributes to the
two real places, which turns the assertion $2^{28}$ at $k=12$ into a theorem. The
class number is a value, returned independently by the kernel's own tower and
Rédei distillations, the latter with $4$-rank one matching $[32,2]$.

The three types were declared \texttt{Type 0} with no inhabitants: placeholders
for constructions \Lean{} does not have. Constructing class fields is not what is
wanted here. Each is recorded instead as the twelve-slot type it occupies, the
narrow field differing from the wide in the parity slot alone, since ramification
at the real places is a parity condition. Nothing is claimed of them beyond
position.

Thirty axioms become none, and the module still carries no \texttt{sorry}. The
result of interest is not the count. Discharging the tail invariant showed it was
false. It asserted that the last two cyclic invariants are $(8,2)$ for
$4\le k\le 15$; the tower table printed in the same module gives $k=4$ as
$[32,8,4,2]$, whose last two are $(4,2)$. Stabilisation begins one level later.
The bound was off by one against data sitting a few lines above it, and as an
axiom it could have stayed there indefinitely, since an axiom is consistent with
its own table by fiat. As a theorem it failed immediately and named the level it
failed at.

That is the argument for spending the effort. Reducing assumption count is
tidiness; what the reduction buys is that assertions which contradict the record
announce themselves.

%──────────────────────────────────────────────────────────────────────────────
\section{What follows}
\label{sec:follows}
%──────────────────────────────────────────────────────────────────────────────

The prediction is checkable without solving the $d=2048$ SIC. Computing the ray
class field at the predicted modulus and confirming that its degree, signature,
and Galois structure are consistent with carrying $d/2$ independent moduli under
the pairing induced by the nontrivial automorphism of $\Fd$ is a computation in
the field alone, with no fiducial required. The flat autocorrelation conditions
on the moduli then serve as an independent check, and the condition that the sum
of squares equals $2/(d+1)$ has already been verified numerically in dimension
eight to high precision from a fiducial obtained with no algebraic input at all.

The equiangularity condition itself is verified where a fiducial exists and is
open where none does, and the two should not be conflated. At $d=12$ the
recovered fiducial has been displaced through the full Weyl--Heisenberg group
and all $d^2-1 = 143$ overlaps satisfy $|\langle\psi,D_{p,q}\psi\rangle|^2 =
1/(d+1)$ to better than $10^{-200}$. At $d=2048$ there is nothing to displace:
$1/(d+1) = 1/2049$ is what a fiducial there would satisfy, and computing that
quotient from $d$ asserts nothing about the existence of a vector attaining it.
The numerical search in that dimension is closed rather than stalled, with its
best objective value far from the floor and a residual flat across every row,
column, and $2$-adic stratum, so no localised obstruction remains to attack;
and the arithmetic route is blocked at the $L$-values of the predicted modulus,
where per-character evaluation exhausts memory and the conductor ascent is
superpolynomial with a climbing exponent. The conductor is what this paper
determines; the fiducial is not, and the distinction is the whole content of
\Cref{sec:follows}.

Dimension twenty no longer needs reconciling. Its coinvariant count falls short
because $5$-torsion is absent from every modulus that keeps the two-part in
range, as \Cref{sec:classgroup} sets out. What it costs is the generality of the
count, not the $d=16$ reading.

\Cref{conj:identification} suggests one further computation, and it is cheap.
The sixteen-power dimensions are $d = 2^{n}$ with $n \equiv 3 \pmod 4$, of which
$d=8$ is already computed and $d=128$ is the first untested. Confirming that the
$d=128$ conductor is $\pt^{8}\infty_1$ with $2$ inert, so that its generator is
$16^{2}$, would put two computed points on the progression instead of one. That
is a ray class computation in a single field and it does not require a fiducial.
A failure there would settle the matter in the other direction, which is the
useful property of the test.

%──────────────────────────────────────────────────────────────────────────────
\section*{Artifact statement}
The field computations were performed in PARI/GP~\cite{PARI}. Everything this
paper claims is frozen on branch \texttt{crystalline/sic-moduli-conductor-2026-08-08}
of \url{https://github.com/umpolungfish/p4rakernel}, under tag
\texttt{crystalline-sic-moduli-conductor-v2}, and that branch is self-contained:
it carries the \Lean{} modules for the tower and the per-dimension report, both
*sans* sorry and the tower module axiom-free, the recovered $d=12$ fiducial that
the equiangularity check runs on, the kernel module that instantiates the
arithmetic of \Cref{sec:unit}, and a table mapping each claim above to the
command that prints it. Running \texttt{./verify\_sic\_moduli.sh} at the root of
that branch elaborates the modules through the \Lean{} kernel and prints,
dimension by dimension across $d = 2, 4, 8, 12, 16, 20, 2048$, what was checked
and the axiom dependency of each headline theorem, so the boundary between what
is proved and what enters as computed field data is visible in the output.

The same branch pins a commit of the mOMonadOS bare-metal kernel
(\url{https://github.com/umpolungfish/momonad_os}), where the arithmetic runs
rather than being described. Booting it and issuing \texttt{d2048 exact} prints
the unit, the two ramified generators with their norms, the exact monomial with
its norm, its rationalized reading, and the four Gauss sums with the radical
each contributes; \texttt{d2048 scaling} prints the measured representation
widths; \texttt{d2048 crossover} prints the threshold under each gate model; and
\texttt{d2048 welch} displaces the $d=12$ fiducial through the full
Weyl--Heisenberg group and reports all $143$ overlaps, together with the fact
that at $d=2048$ there is no vector to displace. Every field computation
reported here is reproducible from the PARI/GP calls named in the text: the
modulus is read off with \texttt{rnfconductor}, and the candidate algebraic
relations with \texttt{algdep}.

%──────────────────────────────────────────────────────────────────────────────
\section*{Acknowledgements}
The author would like to thank Harry T.\ Larson for imparting the importance of catching rising problems and never letting them go~\cite{Larson1961}. The conductor was found by refusing to accept a tool's refusal at face value: when \texttt{bnrstark} declined the coordinate field for having nonzero $r_2$, the question became which field the moduli actually generate — and that question had an answer.

%──────────────────────────────────────────────────────────────────────────────
\begin{thebibliography}{9}
%──────────────────────────────────────────────────────────────────────────────

\bibitem{Zauner1999}
G.\ Zauner, \emph{Quantendesigns: Grundz\"uge einer nichtkommutativen Designtheorie}, Ph.D.\ thesis, Universit\"at Wien, 1999.

\bibitem{Appleby2013}
D.\ M.\ Appleby, H.\ Yadsan-Appleby, and G.\ Zauner,
\emph{Galois automorphisms of a symmetric measurement},
Quantum Inf.\ Comput.\ \textbf{13} (2013), 672--720.

\bibitem{Appleby2017}
M.\ Appleby, S.\ Flammia, G.\ McConnell, and J.\ Yard,
\emph{SICs and algebraic number theory},
Found.\ Phys.\ \textbf{47} (2017), 1042--1059.

\bibitem{AFMY2020}
M.\ Appleby, S.\ Flammia, G.\ McConnell, and J.\ Yard,
\emph{Generating ray class fields of real quadratic fields via complex equiangular lines}, Acta Arith.\ \textbf{192} (2020), 211--233.

\bibitem{Belnap1977}
N.\ D.\ Belnap, \emph{A useful four-valued logic}, in J.\ M.\ Dunn and
G.\ Epstein (eds.), \emph{Modern Uses of Multiple-Valued Logic},
Reidel, 1977, 5--37.

\bibitem{Shramko2005}
Y.\ Shramko and H.\ Wansing, \emph{Some useful 16-valued logics: how a
computer network should think}, J.\ Philos.\ Logic \textbf{34} (2005), 121--153.

\bibitem{PARI}
The PARI~Group, \emph{PARI/GP version 2.13}, Univ.\ Bordeaux, 2021,
\texttt{https://pari.math.u-bordeaux.fr/}.

\bibitem{Larson1961}
H.\ T.\ Larson,
\textit{Catch a Rising Problem\ldots\ and Never Ever Let it Go},
\emph{IRE Transactions on Engineering Management} \textbf{8}(4) (1961), 173--174.
\url{https://ieeexplore.ieee.org/stamp/stamp.jsp?arnumber=1641382}

\end{thebibliography}

\end{document}
